\documentclass[10pt,a4paper]{article}

%double si on veut une version prof et une version élève. simple si on veut une seule version.
\newcommand{\buildmode}{simple}

\InputIfFileExists{version.tex}{}{}
% Version (définie par le script)
\providecommand{\version}{eleve}

\usepackage{feuilleinterro}

%%%à modifier à chaque fois

\renewcommand{\niveau}{2nde}
\renewcommand{\chapitre}{Proportions, pourcentages}
\renewcommand{\numchap}{0.2}
\renewcommand{\theme}{Statistiques}

\begin{document}

% =========================================================
% PAGE 1 : VERSIONS A ET B
% =========================================================

\begin{sujets}

\interroversion{A}

\textbf{Exercice 1 -- Tableau de valeurs \hfill /2}

On considère la fonction \(f\) définie sur \(\mathbb{R}\) par \(f(x)=x^2-2x-3\).

Complétez le tableau de valeurs suivant (une valeur est déjà calculée à titre d'exemple : \(f(1)=1^2-2\times 1-3=-4\)).

\begin{center}
    \renewcommand{\arraystretch}{1.5}
    \setlength{\tabcolsep}{0.25cm}
    \begin{tabular}{|*{7}{c|}}
    \hline
    \(x\) & \(-2\) & \(-1\) & \(0\) & \(1\) & \(2\) & \(3\) \\
    \hline
    \(f(x)\) & \(5\) & & \(-3\) & \(-4\) & & \(0\) \\
    \hline
    \end{tabular}
\end{center}

\vspace{3cm}

\textbf{Exercice 2 -- Pourcentages \hfill /3}

\begin{tasks}(2)
    \task Calculer \(10\%\) de \(90\).
    \task Poser et effectuer le calcul de \(35\%\) de \(80\).
    \task Calculer \(20\%\) de \(15\%\) de \(300\).
    \task Calculer le tiers de \(21\).
\end{tasks}

\columnbreak

\interroversion{B}

\textbf{Exercice 1 -- Tableau de valeurs \hfill /2}

On considère la fonction \(f\) définie sur \(\mathbb{R}\) par \(f(x)=x^2+x-2\).

Complétez le tableau de valeurs suivant (une valeur est déjà calculée à titre d'exemple : \(f(2)=2^2+2-2=4\)).

\begin{center}
    \renewcommand{\arraystretch}{1.5}
    \setlength{\tabcolsep}{0.25cm}
    \begin{tabular}{|*{7}{c|}}
    \hline
    \(x\) & \(-2\) & \(-1\) & \(0\) & \(1\) & \(2\) & \(3\) \\
    \hline
    \(f(x)\) & \(0\) & & \(-2\) & \(0\) & \(4\) &  \\
    \hline
    \end{tabular}
\end{center}

\vspace{3cm}

\textbf{Exercice 2 -- Pourcentages \hfill /3}

\begin{tasks}(2)
    \task Calculer \(20\%\) de \(50\).
    \task Poser et effectuer le calcul de \(24\%\) de \(75\).
    \task Calculer \(10\%\) de \(40\%\) de \(500\).
    \task Calculer le quart de \(28\).
\end{tasks}

\end{sujets}

\newpage

% =========================================================
% PAGE 2 : VERSIONS C ET D
% =========================================================

\begin{sujets}

\interroversion{C}

\textbf{Exercice 1 -- Tableau de valeurs \hfill /2}

On considère la fonction \(f\) définie sur \(\mathbb{R}\) par \(f(x)=x^2-3x+2\).

Complétez le tableau de valeurs suivant (une valeur est déjà calculée à titre d'exemple : \(f(-1)=(-1)^2-3\times(-1)+2=6\)).

\begin{center}
    \renewcommand{\arraystretch}{1.5}
    \setlength{\tabcolsep}{0.25cm}
    \begin{tabular}{|*{7}{c|}}
    \hline
    \(x\) & \(-2\) & \(-1\) & \(0\) & \(1\) & \(2\) & \(3\) \\
    \hline
    \(f(x)\) & & \(6\) & \(2\) & \(0\) & \(0\) & \\
    \hline
    \end{tabular}
\end{center}

\vspace{3cm}

\textbf{Exercice 2 -- Pourcentages \hfill /3}

\begin{tasks}(2)
    \task Calculer \(25\%\) de \(40\).
    \task Poser et effectuer le calcul de \(18\%\) de \(50\).
    \task Calculer \(50\%\) de \(20\%\) de \(200\).
    \task Calculer le tiers de \(18\).
\end{tasks}

\columnbreak

\interroversion{D}

\textbf{Exercice 1 -- Tableau de valeurs \hfill /2}

On considère la fonction \(f\) définie sur \(\mathbb{R}\) par \(f(x)=x^2+2x-1\).

Complétez le tableau de valeurs suivant (une valeur est déjà calculée à titre d'exemple : \(f(0)=0^2+2\times 0-1=-1\)).

\begin{center}
    \renewcommand{\arraystretch}{1.5}
    \setlength{\tabcolsep}{0.25cm}
    \begin{tabular}{|*{7}{c|}}
    \hline
    \(x\) & \(-2\) & \(-1\) & \(0\) & \(1\) & \(2\) & \(3\) \\
    \hline
    \(f(x)\) & \(-1\) & & \(-1\) & \(2\) & & \(14\) \\
    \hline
    \end{tabular}
\end{center}

\vspace{3cm}

\textbf{Exercice 2 -- Pourcentages \hfill /3}

\begin{tasks}(2)
    \task Calculer \(10\%\) de \(70\).
    \task Poser et effectuer le calcul de \(32\%\) de \(75\).
    \task Calculer \(25\%\) de \(40\%\) de \(240\).
    \task Calculer le quart de \(36\).
\end{tasks}

\end{sujets}

\end{document}
